\scanpage{121}
\begin{center}- 93 -\end{center}
\margin{n\textsuperscript{o} 371}

$\bigwedge(N)$, où $N$ est un sous-module de type fini de $M$; il suffit alors d'appliquer le cor.3 de la prop.2 du n\textsuperscript{o} 1.

COROLLAIRE 5.- \uline{Soient} $N$, $P$ \uline{deux sous-modules d'un} $A$-\uline{module} $M$ \uline{tels que} $N\cap P$ \uline{soit facteur direct de} $N$ \uline{et de} $P$ \uline{et} $N+P$ \uline{facteur direct de} $M$. \uline{Alors} $\bigwedge(N)$ \uline{et} $\bigwedge(P)$ \uline{s'identifient canoniquement à des sous-algèbres de} $\bigwedge(M)$, \uline{et} $\bigwedge(N\cap P)$ \uline{à l'intersection} $\bigwedge(N)\cap\bigwedge(P)$.

En effet, si $E$, $F$ sont des supplémentaires de $N\cap P$ dans $N$ et $P$ respectivement et $G$ un supplémentaire de $N+P$ dans $M$, $M$ est somme directe des quatre sous-modules $N\cap P$, $E$, $F$, $G$; donc $\bigwedge(N)$, $\bigwedge(P)$ et $\bigwedge(N\cap P)$ s'identifient à des sous-algèbres de $\bigwedge(M)$ (cor.3). En outre, en vertu du cor.1, $\mathop{\bigwedge}\limits^n N$ (resp. $\mathop{\bigwedge}\limits^n P$) s'identifie à la somme directe des sous-modules $\mathop{\bigwedge}\limits^p(N\cap P)\otimes\mathop{\bigwedge}\limits^{n-p}E$ (resp. $\mathop{\bigwedge}\limits^p(N\cap P)\otimes\mathop{\bigwedge}\limits^{n-p}F$) pour $0\leq p\leq n$; on a donc bien $\bigwedge(N\cap P)=\bigwedge(N)\cap\bigwedge(P)$.

PROPOSITION 8.- \uline{Soit} $(e_\lambda)_{\lambda\in L}$ \uline{une base d'un} $A$-\uline{module} $M$; \uline{munissons} $L$ \uline{d'une structure d'ordre total}, \uline{et pour toute partie finie} $H$ \uline{de} $L$, \uline{posons}
\[
\text{(11)}\qquad e_H=e_{\lambda_1}\wedge e_{\lambda_2}\wedge\cdots\wedge e_{\lambda_n}
\]
\uline{si} $(\lambda_i)_{1\leq i\leq n}$ \uline{est la suite strictement croissante des éléments de} $H$. \uline{Alors les éléments homogènes} $e_H$, \uline{où} $H$ \uline{parcourt l'ensemble} $\unsure{\mathfrak F|\mathrm F}(L)$ \uline{des parties finies de} $L$, \uline{forment une base de} $\bigwedge(M)$.

Supposons d'abord $L$ \uline{fini}, soit par exemple $L=[\![1,m]\!]$ muni de l'ordre usuel; alors la proposition résulte de la prop.7, car si on pose $M_i=Ae_i$ pour $1\leq i\leq m$, $M$ est somme directe des $M_i$ et il résulte aussitôt du cor.3 de la prop.2 du n\textsuperscript{o} 1 que $\bigwedge(M_i)=A\oplus M_i$ (autrement dit $\mathop{\bigwedge}\limits^p M_i=0$ pour $p\geq2$); la conclusion résulte alors du cor.1 de la prop.7 et du chap.II, 3e éd., § 3, n\textsuperscript{o} 4, prop.4.

Supposons ensuite $L$ infini et pour toute partie finie $J$ de $L$, soit $M_J$ le sous-module de $M$ ayant pour base les $e_\lambda$ tels que $\lambda\in J$; il est clair que $M$ est limite
\margin{120}
